\import{algebra/group/generated.tex}
\import{algebra/ring/reducts.tex}
\import{algebra/ring/subring.tex}
\import{algebra/ring/natural.tex}
\import{algebra/ring/binomial.tex}
\import{algebra/magma/subset.tex}
\import{algebra/monoid/subset.tex}
\import{algebra/semigroup/positive-power.tex}
\import{algebra/group/finite-sum.tex}

\section{Additively generated products of subsets}

The product below first multiplies elements pointwise and then takes their
additive span.  It is therefore distinct from the raw pointwise product used
for boundedness questions.

\begin{definition}\label{additive_subset_product}
    $\additiveSubsetProduct{R}{A}{B}=
        \generatedSubgroup{
            \pointwiseProduct{\MultiplicativeMonoid{R}}{A}{B}}
            {\AdditiveGroup{R}}$.
\end{definition}

\begin{proposition}\label{additive_subset_product_type}
    Let $R$ be a ring.
    Suppose $A,B\subseteq\carrier[R]$.
    Then $\additiveSubsetProduct{R}{A}{B}\subseteq\carrier[R]$.
\end{proposition}
\begin{proof}
    Let $G=\AdditiveGroup{R}$ and $M=\MultiplicativeMonoid{R}$.
    $G$ is a group by \obv{additive_group_is_group}.
    $M$ is a monoid by \obv{multiplicative_monoid_is_monoid}.
    $\carrier[G]=\carrier[R]$ and $\carrier[M]=\carrier[R]$ by
        \auto{additive_group_fields,multiplicative_monoid_fields}.
    Hence $A,B\subseteq\carrier[M]$.
    Hence $\pointwiseProduct{M}{A}{B}\subseteq\carrier[M]$ by
        \obv{pointwise_product_type}.
    Hence $\pointwiseProduct{M}{A}{B}\subseteq\carrier[G]$.
    $\generatedSubgroup{\pointwiseProduct{M}{A}{B}}{G}
        \subseteq\carrier[G]$ by
        \obv{generated_subgroup_subseteq_carrier}.
    Follows by \auto{additive_subset_product}.
\end{proof}

\begin{lemma}\label{additive_subgroup_contains_natural_multiple}
    Let $R$ be a ring and $n$ be a natural number.
    Suppose $\subgroup{H}{\AdditiveGroup{R}}$ and $x\in H$.
    Then $\ringMul[R](\ringNatural{R}{n},x)\in H$.
\end{lemma}
\begin{proof}
    Let $G=\AdditiveGroup{R}$.
    $\carrier[G]=\carrier[R]$ and $\mul[G]=\ringAdd[R]$ and
        $\neutral[G]=\ringZero[R]$ by
        \auto{additive_group_fields}.
    $\subgroup{H}{G}$.
    $H\subseteq\carrier[G]$ by \obv{subgroup_subseteq}.
    Hence $x\in\carrier[G]$ by \obv{subseteq}.
    Hence $x\in\carrier[R]$.
    Let $C=\{k\in\naturals\mid
        \ringMul[R](\ringNatural{R}{k},x)\in H\}$.
    $\ringNatural{R}{\emptyset}=\ringZero[R]$ by
        \obv{ring_natural_zero}.
    $\ringMul[R](\ringZero[R],x)=\ringZero[R]$ by
        \obv{ring_mul_zero_left}.
    $\neutral[G]\in H$ by \obv{subgroup_neutral}.
    Hence $\ringMul[R](\ringNatural{R}{\emptyset},x)\in H$.
    Hence $\emptyset\in C$ by \obv{omega_empty}.
    Show that for every $k\in C$ we have $\suc{k}\in C$.
    \begin{subproof}
        Fix $k\in C$.
        Hence $k$ is a natural number and
            $\ringMul[R](\ringNatural{R}{k},x)\in H$.
        $\ringNatural{R}{k},\ringOne[R]\in\carrier[R]$ by
            \obv{ring_natural_type,ring_one_type}.
        $\ringNatural{R}{\suc{k}}
            =\ringAdd[R](\ringNatural{R}{k},\ringOne[R])$ by
            \obv{ring_natural_successor}.
        $\ringMul[R](\ringNatural{R}{\suc{k}},x)
            =\ringAdd[R](\ringMul[R](\ringNatural{R}{k},x),x)$ by
            \auto{ring_distributive_right,ring_mul_one_left}.
        $\mul[G](\ringMul[R](\ringNatural{R}{k},x),x)\in H$ by
            \obv{subgroup_mul}.
        Hence $\ringMul[R](\ringNatural{R}{\suc{k}},x)\in H$.
        $\suc{k}$ is a natural number by \obv{omega_successor}.
    \end{subproof}
    Hence $C$ is an inductive set by \obv{inductive_set}.
    Hence every natural number is an element of $C$ by
        \obv{omega_minimal}.
    Hence $n\in C$.
\end{proof}

\begin{theorem}\label{additive_subset_product_is_subgroup}
    Let $R$ be a ring.
    Suppose $A,B\subseteq\carrier[R]$.
    Then
        $\subgroup{\additiveSubsetProduct{R}{A}{B}}
            {\AdditiveGroup{R}}$.
\end{theorem}
\begin{proof}
    Let $G=\AdditiveGroup{R}$ and $M=\MultiplicativeMonoid{R}$.
    $G$ is a group by \obv{additive_group_is_group}.
    $M$ is a monoid by \obv{multiplicative_monoid_is_monoid}.
    $\carrier[G]=\carrier[R]$ and $\carrier[M]=\carrier[R]$ by
        \auto{additive_group_fields,multiplicative_monoid_fields}.
    Hence $A,B\subseteq\carrier[M]$.
    Hence $\pointwiseProduct{M}{A}{B}\subseteq\carrier[M]$ by
        \obv{pointwise_product_type}.
    Hence $\pointwiseProduct{M}{A}{B}\subseteq\carrier[G]$.
    $\subgroup{\generatedSubgroup{
        \pointwiseProduct{M}{A}{B}}{G}}{G}$ by
        \obv{generated_subgroup_is_subgroup}.
    Follows by \auto{additive_subset_product}.
\end{proof}

\begin{proposition}\label{pointwise_product_subseteq_additive_subset_product}
    Let $R$ be a ring.
    Suppose $A,B\subseteq\carrier[R]$.
    Then $\pointwiseProduct{\MultiplicativeMonoid{R}}{A}{B}
        \subseteq\additiveSubsetProduct{R}{A}{B}$.
\end{proposition}
\begin{proof}
    Let $G=\AdditiveGroup{R}$ and $M=\MultiplicativeMonoid{R}$.
    $G$ is a group by \obv{additive_group_is_group}.
    $M$ is a monoid by \obv{multiplicative_monoid_is_monoid}.
    Hence $M$ is a magma by \obv{monoid_is_magma}.
    $\carrier[G]=\carrier[R]$ and $\carrier[M]=\carrier[R]$ by
        \auto{additive_group_fields,multiplicative_monoid_fields}.
    Hence $A,B\subseteq\carrier[M]$.
    Hence $\pointwiseProduct{M}{A}{B}\subseteq\carrier[G]$ by
        \obv{pointwise_product_type}.
    $\pointwiseProduct{M}{A}{B}\subseteq
        \generatedSubgroup{\pointwiseProduct{M}{A}{B}}{G}$ by
        \obv{generated_subgroup_contains_generators}.
    Follows by \auto{additive_subset_product}.
\end{proof}

\begin{theorem}\label{ring_power_of_sum_in_power_envelope_product}
    Let $R$ be a commutative ring and $x,y\in\carrier[R]$.
    Let $n$ be a natural number.
    Then
        $\ringPower{R}{\ringAdd[R](x,y)}{n}\in
            \additiveSubsetProduct{R}{
                \pointwiseSubsetPowers{\MultiplicativeMonoid{R}}{\{x\}}}{
                \pointwiseSubsetPowers{\MultiplicativeMonoid{R}}{\{y\}}}$.
\end{theorem}
\begin{proof}
    Let $G=\AdditiveGroup{R}$ and $M=\MultiplicativeMonoid{R}$.
    Let $X=\pointwiseSubsetPowers{M}{\{x\}}$ and
        $Y=\pointwiseSubsetPowers{M}{\{y\}}$.
    Let $H=\additiveSubsetProduct{R}{X}{Y}$.
    $R$ is a ring by \obv{commutative_ring_is_ring}.
    $G$ is a commutative group by
        \obv{additive_group_is_exact_commutative_group}.
    $M$ is a monoid by \obv{multiplicative_monoid_is_exact_monoid}.
    $\carrier[G]=\carrier[R]$ and $\carrier[M]=\carrier[R]$ and
        $\mul[M]=\ringMul[R]$ by
        \auto{additive_group_fields,multiplicative_monoid_fields}.
    Hence $x,y\in\carrier[M]$.
    $\{x\},\{y\}\subseteq\carrier[M]$ by
        \auto{subseteq,singleton_iff}.
    Hence $\pointwiseSubsetPowers{M}{\{x\}},
        \pointwiseSubsetPowers{M}{\{y\}}\subseteq\carrier[M]$ by
        \obv{pointwise_subset_powers_type}.
    Hence $X,Y\subseteq\carrier[M]$.
    Hence $X,Y\subseteq\carrier[R]$.
    $\subgroup{\additiveSubsetProduct{R}{X}{Y}}{\AdditiveGroup{R}}$ by
        \obv{additive_subset_product_is_subgroup}.
    Hence $\subgroup{H}{G}$.
    Let $f=\Apply{\graphOnOp}{
        \suc{n},\carrier[R],
        \fun{k:\Set}{
            \ringMul[R](
                \ringNatural{R}{\naturalChoose{n}{k}},
                \ringMul[R](
                    \ringPower{R}{x}{\naturalDifference{n}{k}},
                    \ringPower{R}{y}{k}))}}$.
    $f\in\funs{\suc{n}}{\carrier[R]}$ by
        \auto{ring_binomial_summand_graph_type}.
    Hence $f\in\funs{\suc{n}}{\carrier[G]}$ by
        \obv{additive_group_fields}.
    Show that for every $k\in\suc{n}$ we have $f(k)\in H$.
    \begin{subproof}
        Fix $k\in\suc{n}$.
        $k$ is a natural number by
            \auto{omega_successor,naturals_is_transitive_set,
                transitiveset_iff_subseteq,subseteq}.
        $k$ and $n$ are ordinals by \obv{natural_number_is_ordinal}.
        $k=n$ or $k\in n$ by \obv{suc_elim}.
        Hence $k\subseteq n$ by
            \obv{subseteq_refl,ordinal_elem_implies_subseteq}.
        $\naturalDifference{n}{k}$ is a natural number by
            \auto{natural_difference_specification}.
        $\naturalChoose{n}{k}$ is a natural number by
            \obv{natural_choose_type}.
        $\ringPower{R}{x}{\naturalDifference{n}{k}}
            =\monoidPower{M}{x}{\naturalDifference{n}{k}}$ and
            $\ringPower{R}{y}{k}=\monoidPower{M}{y}{k}$ by
            \obv{ring_power}.
        $\pointwiseSubsetPower{M}{\{x\}}{\naturalDifference{n}{k}}
            =\{\monoidPower{M}{x}{\naturalDifference{n}{k}}\}$ and
            $\pointwiseSubsetPower{M}{\{y\}}{k}
            =\{\monoidPower{M}{y}{k}\}$ by
            \obv{pointwise_subset_power_singleton}.
        Hence $\ringPower{R}{x}{\naturalDifference{n}{k}}\in X$ and
            $\ringPower{R}{y}{k}\in Y$ by
            \auto{singleton_iff,pointwise_subset_powers_iff}.
        Hence $\ringMul[R](
                \ringPower{R}{x}{\naturalDifference{n}{k}},
                \ringPower{R}{y}{k})
            \in\pointwiseProduct{M}{X}{Y}$ by
            \auto{pointwise_product_iff}.
        $\pointwiseProduct{\MultiplicativeMonoid{R}}{X}{Y}
            \subseteq\additiveSubsetProduct{R}{X}{Y}$ by
            \obv{pointwise_product_subseteq_additive_subset_product}.
        Hence $\pointwiseProduct{M}{X}{Y}
            \subseteq\additiveSubsetProduct{R}{X}{Y}$.
        Hence $\ringMul[R](
                \ringPower{R}{x}{\naturalDifference{n}{k}},
                \ringPower{R}{y}{k})
            \in\additiveSubsetProduct{R}{X}{Y}$ by \obv{subseteq}.
        Hence $\ringMul[R](
                \ringPower{R}{x}{\naturalDifference{n}{k}},
                \ringPower{R}{y}{k})\in H$.
        $\ringMul[R](
                \ringNatural{R}{\naturalChoose{n}{k}},
                \ringMul[R](
                    \ringPower{R}{x}{\naturalDifference{n}{k}},
                    \ringPower{R}{y}{k}))
            \in\additiveSubsetProduct{R}{X}{Y}$ by
            \obv{additive_subgroup_contains_natural_multiple}.
        Hence $\ringMul[R](
                \ringNatural{R}{\naturalChoose{n}{k}},
                \ringMul[R](
                    \ringPower{R}{x}{\naturalDifference{n}{k}},
                    \ringPower{R}{y}{k}))\in H$.
        $f(k)=\ringMul[R](
                \ringNatural{R}{\naturalChoose{n}{k}},
                \ringMul[R](
                    \ringPower{R}{x}{\naturalDifference{n}{k}},
                    \ringPower{R}{y}{k}))$ by
            \auto{logical_graph_on_apply}.
    \end{subproof}
    $\suc{n}$ is finite by
        \obv{omega_successor,natural_number_is_finite}.
    Hence $\finiteSum{G}{\suc{n}}{f}\in H$ by
        \obv{subgroup_contains_finite_sum}.
    $\ringBinomialExpansion{R}{x}{y}{n}
        =\finiteSum{G}{\suc{n}}{f}$ by
        \obv{ring_binomial_expansion}.
    $\ringPower{R}{\ringAdd[R](x,y)}{n}
        =\ringBinomialExpansion{R}{x}{y}{n}$ by
        \obv{ring_binomial_theorem}.
    Hence $\ringPower{R}{\ringAdd[R](x,y)}{n}\in H$.
\end{proof}

\begin{proposition}\label{additive_subset_product_monotone}
    Let $R$ be a ring.
    Suppose $A,A',B,B'\subseteq\carrier[R]$.
    Suppose $A\subseteq A'$ and $B\subseteq B'$.
    Then $\additiveSubsetProduct{R}{A}{B}
        \subseteq\additiveSubsetProduct{R}{A'}{B'}$.
\end{proposition}
\begin{proof}
    Let $G=\AdditiveGroup{R}$ and $M=\MultiplicativeMonoid{R}$.
    $G$ is a group by \obv{additive_group_is_group}.
    $M$ is a monoid by \obv{multiplicative_monoid_is_monoid}.
    Hence $M$ is a magma by \obv{monoid_is_magma}.
    $\carrier[G]=\carrier[R]$ and $\carrier[M]=\carrier[R]$ by
        \auto{additive_group_fields,multiplicative_monoid_fields}.
    Hence $A,A',B,B'\subseteq\carrier[M]$.
    $\pointwiseProduct{M}{A}{B}
        \subseteq\pointwiseProduct{M}{A'}{B'}$ by
        \obv{pointwise_product_monotone}.
    $\pointwiseProduct{M}{A}{B},
        \pointwiseProduct{M}{A'}{B'}\subseteq\carrier[G]$ by
        \obv{pointwise_product_type}.
    Follows by
        \auto{additive_subset_product,generated_subgroup_monotone}.
\end{proof}

\subsection{Multiplication and additive generation}

\begin{lemma}\label{generated_subgroup_closed_under_right_ring_multiplication}
    Let $R$ be a ring and $G=\AdditiveGroup{R}$.
    Suppose $A\subseteq\carrier[R]$ and $b\in\carrier[R]$.
    Suppose $\subgroup{H}{G}$.
    Suppose for every $a\in A$ we have $\ringMul[R](a,b)\in H$.
    Then for every $x\in\generatedSubgroup{A}{G}$ we have
        $\ringMul[R](x,b)\in H$.
\end{lemma}
\begin{proof}
    $\carrier[G]=\carrier[R]$ and
        $\mul[G]=\ringAdd[R]$ and
        $\neutral[G]=\ringZero[R]$ and
        $\groupInv[G]=\ringNeg[R]$ by
        \auto{additive_group_fields}.
    Let $K=\{x\in\carrier[R]\mid \ringMul[R](x,b)\in H\}$.
    $K\subseteq\carrier[G]$ by \obv{subseteq}.
    $\neutral[G]\in H$ by \obv{subgroup_neutral}.
    $\ringMul[R](\ringZero[R],b)=\ringZero[R]$ by
        \obv{ring_mul_zero_left}.
    $\neutral[G]\in\carrier[R]$ by \obv{ring_zero_type}.
    $\ringMul[R](\neutral[G],b)\in H$.
    Hence $\neutral[G]\in K$.
    Show that for all $x,y\in K$ we have $\mul[G](x,y)\in K$.
    \begin{subproof}
        Fix $x,y\in K$.
        Hence $x,y\in\carrier[R]$ and
            $\ringMul[R](x,b),\ringMul[R](y,b)\in H$.
        $\mul[G](\ringMul[R](x,b),\ringMul[R](y,b))\in H$ by
            \obv{subgroup_mul}.
        Hence $\ringAdd[R](\ringMul[R](x,b),\ringMul[R](y,b))\in H$.
        $\ringMul[R](\ringAdd[R](x,y),b)
            =\ringAdd[R](\ringMul[R](x,b),\ringMul[R](y,b))$ by
            \obv{ring_distributive_right}.
        $\ringAdd[R](x,y)\in\carrier[R]$ by \obv{ring_add_type}.
        Follows by \auto.
    \end{subproof}
    Show that for every $x\in K$ we have $\groupInv[G](x)\in K$.
    \begin{subproof}
        Fix $x\in K$.
        Hence $x\in\carrier[R]$ and $\ringMul[R](x,b)\in H$.
        $\groupInv[G](\ringMul[R](x,b))\in H$ by
            \obv{subgroup_inv}.
        Hence $\ringNeg[R](\ringMul[R](x,b))\in H$.
        $\ringMul[R](\ringNeg[R](x),b)
            =\ringNeg[R](\ringMul[R](x,b))$ by
            \obv{ring_neg_product_left}.
        $\ringNeg[R](x)\in\carrier[R]$ by \obv{ring_neg_type}.
    \end{subproof}
    Hence $\subgroup{K}{G}$ by \obv{subgroup}.
    Show that $A\subseteq K$.
    \begin{subproof}
        Follows by \obv{subseteq}.
    \end{subproof}
    Hence $\generatedSubgroup{A}{G}\subseteq K$ by
        \obv{generated_subgroup_least}.
    Follows by \auto{subseteq}.
\end{proof}

\begin{lemma}\label{generated_subgroup_closed_under_left_ring_multiplication}
    Let $R$ be a ring and $G=\AdditiveGroup{R}$.
    Suppose $A\subseteq\carrier[R]$ and $a\in\carrier[R]$.
    Suppose $\subgroup{H}{G}$.
    Suppose for every $b\in A$ we have $\ringMul[R](a,b)\in H$.
    Then for every $y\in\generatedSubgroup{A}{G}$ we have
        $\ringMul[R](a,y)\in H$.
\end{lemma}
\begin{proof}
    $\carrier[G]=\carrier[R]$ and
        $\mul[G]=\ringAdd[R]$ and
        $\neutral[G]=\ringZero[R]$ and
        $\groupInv[G]=\ringNeg[R]$ by
        \auto{additive_group_fields}.
    Let $K=\{y\in\carrier[R]\mid \ringMul[R](a,y)\in H\}$.
    $K\subseteq\carrier[G]$ by \obv{subseteq}.
    $\neutral[G]\in H$ by \obv{subgroup_neutral}.
    $\ringMul[R](a,\ringZero[R])=\ringZero[R]$ by
        \obv{ring_mul_zero_right}.
    $\neutral[G]\in\carrier[R]$ by \obv{ring_zero_type}.
    $\ringMul[R](a,\neutral[G])\in H$.
    Hence $\neutral[G]\in K$.
    Show that for all $x,y\in K$ we have $\mul[G](x,y)\in K$.
    \begin{subproof}
        Fix $x,y\in K$.
        Hence $x,y\in\carrier[R]$ and
            $\ringMul[R](a,x),\ringMul[R](a,y)\in H$.
        $\mul[G](\ringMul[R](a,x),\ringMul[R](a,y))\in H$ by
            \obv{subgroup_mul}.
        Hence $\ringAdd[R](\ringMul[R](a,x),\ringMul[R](a,y))\in H$.
        $\ringMul[R](a,\ringAdd[R](x,y))
            =\ringAdd[R](\ringMul[R](a,x),\ringMul[R](a,y))$ by
            \obv{ring_distributive_left}.
        $\ringAdd[R](x,y)\in\carrier[R]$ by \obv{ring_add_type}.
        Follows by \auto.
    \end{subproof}
    Show that for every $y\in K$ we have $\groupInv[G](y)\in K$.
    \begin{subproof}
        Fix $y\in K$.
        Hence $y\in\carrier[R]$ and $\ringMul[R](a,y)\in H$.
        $\groupInv[G](\ringMul[R](a,y))\in H$ by
            \obv{subgroup_inv}.
        Hence $\ringNeg[R](\ringMul[R](a,y))\in H$.
        $\ringMul[R](a,\ringNeg[R](y))
            =\ringNeg[R](\ringMul[R](a,y))$ by
            \obv{ring_neg_product_right}.
        $\ringNeg[R](y)\in\carrier[R]$ by \obv{ring_neg_type}.
    \end{subproof}
    Hence $\subgroup{K}{G}$ by \obv{subgroup}.
    Show that $A\subseteq K$.
    \begin{subproof}
        Follows by \obv{subseteq}.
    \end{subproof}
    Hence $\generatedSubgroup{A}{G}\subseteq K$ by
        \obv{generated_subgroup_least}.
    Follows by \auto{subseteq}.
\end{proof}

\begin{theorem}\label{additive_subset_product_generated_left}
    Let $R$ be a ring and $G=\AdditiveGroup{R}$.
    Suppose $A,B\subseteq\carrier[R]$.
    Then
        $\additiveSubsetProduct{R}{\generatedSubgroup{A}{G}}{B}
            =\additiveSubsetProduct{R}{A}{B}$.
\end{theorem}
\begin{proof}
    Let $M=\MultiplicativeMonoid{R}$ and
        $H=\additiveSubsetProduct{R}{A}{B}$.
    $G$ is a group by \obv{additive_group_is_group}.
    $M$ is a monoid by \obv{multiplicative_monoid_is_monoid}.
    Hence $M$ is a magma by \obv{monoid_is_magma}.
    $\carrier[G]=\carrier[R]$ and $\carrier[M]=\carrier[R]$ and
        $\mul[M]=\ringMul[R]$ by
        \auto{additive_group_fields,multiplicative_monoid_fields}.
    $\generatedSubgroup{A}{G}\subseteq\carrier[R]$ by
        \obv{generated_subgroup_subseteq_carrier}.
    $A\subseteq\carrier[G]$.
    $A\subseteq\generatedSubgroup{A}{G}$ by
        \obv{generated_subgroup_contains_generators}.
    Hence $H\subseteq
        \additiveSubsetProduct{R}{\generatedSubgroup{A}{G}}{B}$ by
        \auto{additive_subset_product_monotone,subseteq_refl}.

    $\subgroup{\additiveSubsetProduct{R}{A}{B}}{\AdditiveGroup{R}}$ by
        \obv{additive_subset_product_is_subgroup}.
    Hence $\subgroup{H}{G}$.
    Show that
        $\pointwiseProduct{M}{\generatedSubgroup{A}{G}}{B}\subseteq H$.
    \begin{subproof}
        Show that every element of
            $\pointwiseProduct{M}{\generatedSubgroup{A}{G}}{B}$ is an
            element of $H$.
        \begin{subproof}
            Fix $z\in
                \pointwiseProduct{M}{\generatedSubgroup{A}{G}}{B}$.
            Take $x,b$ such that $x\in\generatedSubgroup{A}{G}$ and
                $b\in B$ and $z=\mul[M](x,b)$ by
                \obv{pointwise_product_iff}.
            $b\in\carrier[R]$ by \obv{subseteq}.
            Show that for every $a\in A$ we have $\ringMul[R](a,b)\in H$.
            \begin{subproof}
                Fix $a\in A$.
                $\mul[M](a,b)\in\pointwiseProduct{M}{A}{B}$ by
                    \auto{pointwise_product_iff}.
                Hence $\ringMul[R](a,b)\in
                    \pointwiseProduct{\MultiplicativeMonoid{R}}{A}{B}$.
                $\pointwiseProduct{\MultiplicativeMonoid{R}}{A}{B}
                    \subseteq\additiveSubsetProduct{R}{A}{B}$ by
                    \obv{pointwise_product_subseteq_additive_subset_product}.
                Hence $\ringMul[R](a,b)\in
                    \additiveSubsetProduct{R}{A}{B}$ by
                    \obv{subseteq}.
            \end{subproof}
            Hence $\ringMul[R](x,b)\in
                \additiveSubsetProduct{R}{A}{B}$ by
                \auto{generated_subgroup_closed_under_right_ring_multiplication}.
        \end{subproof}
        Follows by \obv{subseteq}.
    \end{subproof}
    $\pointwiseProduct{M}{\generatedSubgroup{A}{G}}{B}
        \subseteq\carrier[G]$ by \obv{pointwise_product_type}.
    Hence $\generatedSubgroup{
        \pointwiseProduct{M}{\generatedSubgroup{A}{G}}{B}}{G}
        \subseteq H$ by \obv{generated_subgroup_least}.
    Hence $\additiveSubsetProduct{R}{\generatedSubgroup{A}{G}}{B}
        \subseteq H$ by \auto{additive_subset_product}.
    Hence $\additiveSubsetProduct{R}{\generatedSubgroup{A}{G}}{B}=H$ by
        \obv{subseteq_antisymmetric}.
\end{proof}

\begin{theorem}\label{additive_subset_product_generated_right}
    Let $R$ be a ring and $G=\AdditiveGroup{R}$.
    Suppose $A,B\subseteq\carrier[R]$.
    Then
        $\additiveSubsetProduct{R}{A}{\generatedSubgroup{B}{G}}
            =\additiveSubsetProduct{R}{A}{B}$.
\end{theorem}
\begin{proof}
    Let $M=\MultiplicativeMonoid{R}$ and
        $H=\additiveSubsetProduct{R}{A}{B}$.
    $G$ is a group by \obv{additive_group_is_group}.
    $M$ is a monoid by \obv{multiplicative_monoid_is_monoid}.
    Hence $M$ is a magma by \obv{monoid_is_magma}.
    $\carrier[G]=\carrier[R]$ and $\carrier[M]=\carrier[R]$ and
        $\mul[M]=\ringMul[R]$ by
        \auto{additive_group_fields,multiplicative_monoid_fields}.
    $\generatedSubgroup{B}{G}\subseteq\carrier[R]$ by
        \obv{generated_subgroup_subseteq_carrier}.
    $B\subseteq\carrier[G]$.
    $B\subseteq\generatedSubgroup{B}{G}$ by
        \obv{generated_subgroup_contains_generators}.
    Hence $H\subseteq
        \additiveSubsetProduct{R}{A}{\generatedSubgroup{B}{G}}$ by
        \auto{additive_subset_product_monotone,subseteq_refl}.

    $\subgroup{\additiveSubsetProduct{R}{A}{B}}{\AdditiveGroup{R}}$ by
        \obv{additive_subset_product_is_subgroup}.
    Hence $\subgroup{H}{G}$.
    Show that
        $\pointwiseProduct{M}{A}{\generatedSubgroup{B}{G}}\subseteq H$.
    \begin{subproof}
        Show that every element of
            $\pointwiseProduct{M}{A}{\generatedSubgroup{B}{G}}$ is an
            element of $H$.
        \begin{subproof}
            Fix $z\in
                \pointwiseProduct{M}{A}{\generatedSubgroup{B}{G}}$.
            Take $a,y$ such that $a\in A$ and
                $y\in\generatedSubgroup{B}{G}$ and $z=\mul[M](a,y)$ by
                \obv{pointwise_product_iff}.
            $a\in\carrier[R]$ by \obv{subseteq}.
            Show that for every $b\in B$ we have $\ringMul[R](a,b)\in H$.
            \begin{subproof}
                Fix $b\in B$.
                $\mul[M](a,b)\in\pointwiseProduct{M}{A}{B}$ by
                    \obv{pointwise_product_iff}.
                Hence $\ringMul[R](a,b)\in
                    \pointwiseProduct{\MultiplicativeMonoid{R}}{A}{B}$.
                $\pointwiseProduct{\MultiplicativeMonoid{R}}{A}{B}
                    \subseteq\additiveSubsetProduct{R}{A}{B}$ by
                    \obv{pointwise_product_subseteq_additive_subset_product}.
                Hence $\ringMul[R](a,b)\in
                    \additiveSubsetProduct{R}{A}{B}$ by
                    \obv{subseteq}.
            \end{subproof}
            Hence $\ringMul[R](a,y)\in
                \additiveSubsetProduct{R}{A}{B}$ by
                \auto{generated_subgroup_closed_under_left_ring_multiplication}.
        \end{subproof}
        Follows by \obv{subseteq}.
    \end{subproof}
    $\pointwiseProduct{M}{A}{\generatedSubgroup{B}{G}}
        \subseteq\carrier[G]$ by \obv{pointwise_product_type}.
    Hence $\generatedSubgroup{
        \pointwiseProduct{M}{A}{\generatedSubgroup{B}{G}}}{G}
        \subseteq H$ by \obv{generated_subgroup_least}.
    Hence $\additiveSubsetProduct{R}{A}{\generatedSubgroup{B}{G}}
        \subseteq H$ by \auto{additive_subset_product}.
    Hence $\additiveSubsetProduct{R}{A}{\generatedSubgroup{B}{G}}=H$ by
        \obv{subseteq_antisymmetric}.
\end{proof}

\begin{theorem}\label{additive_subset_product_associative}
    Let $R$ be a ring.
    Suppose $A,B,C\subseteq\carrier[R]$.
    Then
        $\additiveSubsetProduct{R}{
            \additiveSubsetProduct{R}{A}{B}}{C}
        =\additiveSubsetProduct{R}{A}{
            \additiveSubsetProduct{R}{B}{C}}$.
\end{theorem}
\begin{proof}
    Let $G=\AdditiveGroup{R}$ and $M=\MultiplicativeMonoid{R}$.
    Let $X=\pointwiseProduct{M}{A}{B}$ and
        $Y=\pointwiseProduct{M}{B}{C}$.
    $M$ is a monoid by \obv{multiplicative_monoid_is_monoid}.
    Hence $M$ is a semigroup by \obv{monoid_implies_semigroup}.
    Hence $M$ is a magma by \obv{semigroup}.
    $\carrier[G]=\carrier[R]$ and $\carrier[M]=\carrier[R]$ by
        \auto{additive_group_fields,multiplicative_monoid_fields}.
    Hence $A,B,C\subseteq\carrier[M]$.
    Hence $X,Y\subseteq\carrier[M]$ by
        \auto{pointwise_product_type}.
    Hence $X,Y\subseteq\carrier[R]$ by \auto.
    $\additiveSubsetProduct{R}{A}{B}=\generatedSubgroup{X}{G}$ and
        $\additiveSubsetProduct{R}{B}{C}=\generatedSubgroup{Y}{G}$ by
        \auto{additive_subset_product}.
    $\additiveSubsetProduct{R}{\generatedSubgroup{X}{G}}{C}
        =\additiveSubsetProduct{R}{X}{C}$ by
        \obv{additive_subset_product_generated_left}.
    $\additiveSubsetProduct{R}{A}{\generatedSubgroup{Y}{G}}
        =\additiveSubsetProduct{R}{A}{Y}$ by
        \obv{additive_subset_product_generated_right}.
    $\pointwiseProduct{M}{X}{C}=\pointwiseProduct{M}{A}{Y}$ by
        \auto{pointwise_product_associative}.
    $\additiveSubsetProduct{R}{X}{C}
        =\generatedSubgroup{\pointwiseProduct{M}{X}{C}}{G}$ and
        $\additiveSubsetProduct{R}{A}{Y}
        =\generatedSubgroup{\pointwiseProduct{M}{A}{Y}}{G}$ by
        \auto{additive_subset_product}.
    Follows by \auto.
\end{proof}

\begin{theorem}\label{commutative_ring_pointwise_product_commutative}
    Let $R$ be a commutative ring.
    Suppose $A,B\subseteq\carrier[R]$.
    Then $\pointwiseProduct{\MultiplicativeMonoid{R}}{A}{B}
        =\pointwiseProduct{\MultiplicativeMonoid{R}}{B}{A}$.
\end{theorem}
\begin{proof}
    Let $M=\MultiplicativeMonoid{R}$.
    $\mul[M]=\ringMul[R]$ by \auto{multiplicative_monoid_fields}.
    Show that $\pointwiseProduct{M}{A}{B}
        =\pointwiseProduct{M}{B}{A}$.
    \begin{subproof}
        Show that every element of $\pointwiseProduct{M}{A}{B}$ is an
            element of $\pointwiseProduct{M}{B}{A}$.
        \begin{subproof}
            Fix $z\in\pointwiseProduct{M}{A}{B}$.
            Take $a,b$ such that $a\in A$ and $b\in B$ and
                $z=\mul[M](a,b)$ by \auto{pointwise_product_iff}.
            $a,b\in\carrier[R]$ by \obv{subseteq}.
            $\ringMul[R](a,b)=\ringMul[R](b,a)$ by
                \obv{commutative_ring_mul_commutative}.
            Hence $z=\mul[M](b,a)$ by \auto.
            Follows by \auto{pointwise_product_iff}.
        \end{subproof}
        Show that every element of $\pointwiseProduct{M}{B}{A}$ is an
            element of $\pointwiseProduct{M}{A}{B}$.
        \begin{subproof}
            Fix $z\in\pointwiseProduct{M}{B}{A}$.
            Take $b,a$ such that $b\in B$ and $a\in A$ and
                $z=\mul[M](b,a)$ by \auto{pointwise_product_iff}.
            $a,b\in\carrier[R]$ by \obv{subseteq}.
            $\ringMul[R](b,a)=\ringMul[R](a,b)$ by
                \obv{commutative_ring_mul_commutative}.
            Hence $z=\mul[M](a,b)$ by \auto.
            Follows by \auto{pointwise_product_iff}.
        \end{subproof}
        Follows by set extensionality.
    \end{subproof}
\end{proof}

\begin{theorem}\label{additive_subset_product_commutative}
    Let $R$ be a commutative ring.
    Suppose $A,B\subseteq\carrier[R]$.
    Then $\additiveSubsetProduct{R}{A}{B}
        =\additiveSubsetProduct{R}{B}{A}$.
\end{theorem}
\begin{proof}
    $\pointwiseProduct{\MultiplicativeMonoid{R}}{A}{B}
        =\pointwiseProduct{\MultiplicativeMonoid{R}}{B}{A}$ by
        \obv{commutative_ring_pointwise_product_commutative}.
    Follows by \auto{additive_subset_product}.
\end{proof}

\subsection{Scalar multiples of subrings}

\begin{theorem}\label{left_scalar_multiple_is_additive_subgroup}
    Let $R$ be a ring.
    Suppose $\subgroup{A}{\AdditiveGroup{R}}$.
    Let $a\in\carrier[R]$.
    Then
        $\subgroup{
            \pointwiseProduct{\MultiplicativeMonoid{R}}{\{a\}}{A}}
            {\AdditiveGroup{R}}$.
\end{theorem}
\begin{proof}
    Let $G=\AdditiveGroup{R}$ and $M=\MultiplicativeMonoid{R}$.
    Let $K=\pointwiseProduct{M}{\{a\}}{A}$.
    $M$ is a monoid by \obv{multiplicative_monoid_is_exact_monoid}.
    Hence $M$ is a magma by \obv{monoid_is_magma}.
    $\carrier[G]=\carrier[R]$ and $\carrier[M]=\carrier[R]$ and
        $\mul[G]=\ringAdd[R]$ and $\neutral[G]=\ringZero[R]$ and
        $\groupInv[G]=\ringNeg[R]$ and $\mul[M]=\ringMul[R]$ by
        \auto{additive_group_fields,multiplicative_monoid_fields}.
    $A\subseteq\carrier[\AdditiveGroup{R}]$ by
        \obv{subgroup_subseteq}.
    Hence $A\subseteq\carrier[G]$ by \auto.
    Hence $A\subseteq\carrier[R]$ by \auto.
    Hence $\{a\},A\subseteq\carrier[M]$ by
        \auto{subseteq,singleton_iff}.
    Hence $K\subseteq\carrier[M]$ by
        \auto{pointwise_product_type}.
    Hence $K\subseteq\carrier[G]$.

    $\ringZero[R]\in A$ by
        \auto{subgroup_neutral,additive_group_fields}.
    $\ringMul[R](a,\ringZero[R])=\ringZero[R]$ by
        \obv{ring_mul_zero_right}.
    Hence $\neutral[G]\in K$ by
        \auto{pointwise_product_iff,singleton_iff}.

    Show that for all $x,y\in K$ we have $\mul[G](x,y)\in K$.
    \begin{subproof}
        Fix $x,y\in K$.
        Take $u\in A$ such that $x=\ringMul[R](a,u)$ by
            \auto{pointwise_product_iff,singleton_iff}.
        Take $v\in A$ such that $y=\ringMul[R](a,v)$ by
            \auto{pointwise_product_iff,singleton_iff}.
        $u,v\in\carrier[R]$ by \obv{subseteq}.
        $\ringAdd[R](u,v)\in A$ by
            \auto{subgroup_mul,additive_group_fields}.
        $\ringAdd[R](x,y)=
            \ringMul[R](a,\ringAdd[R](u,v))$ by
            \auto{ring_distributive_left}.
        Follows by
            \auto{pointwise_product_iff,singleton_iff}.
    \end{subproof}

    Show that for every $x\in K$ we have $\groupInv[G](x)\in K$.
    \begin{subproof}
        Fix $x\in K$.
        Take $u\in A$ such that $x=\ringMul[R](a,u)$ by
            \auto{pointwise_product_iff,singleton_iff}.
        $u\in\carrier[R]$ by \obv{subseteq}.
        $\ringNeg[R](u)\in A$ by
            \auto{subgroup_inv,additive_group_fields}.
        $\ringNeg[R](x)=\ringMul[R](a,\ringNeg[R](u))$ by
            \auto{ring_neg_product_right}.
        Follows by
            \auto{pointwise_product_iff,singleton_iff}.
    \end{subproof}
    Follows by \auto{subgroup}.
\end{proof}

\begin{proposition}\label{additive_subset_product_singleton_subgroup}
    Let $R$ be a ring.
    Suppose $\subgroup{A}{\AdditiveGroup{R}}$.
    Let $a\in\carrier[R]$.
    Then $\additiveSubsetProduct{R}{\{a\}}{A}
        =\pointwiseProduct{\MultiplicativeMonoid{R}}{\{a\}}{A}$.
\end{proposition}
\begin{proof}
    Let $G=\AdditiveGroup{R}$ and $M=\MultiplicativeMonoid{R}$.
    $\subgroup{
        \pointwiseProduct{\MultiplicativeMonoid{R}}{\{a\}}{A}}{
        \AdditiveGroup{R}}$ by
        \obv{left_scalar_multiple_is_additive_subgroup}.
    Hence $\subgroup{\pointwiseProduct{M}{\{a\}}{A}}{G}$ by
        \auto.
    $\generatedSubgroup{\pointwiseProduct{M}{\{a\}}{A}}{G}
        =\pointwiseProduct{M}{\{a\}}{A}$ by
        \obv{generated_subgroup_of_subgroup}.
    Follows by \auto{additive_subset_product}.
\end{proof}

\begin{theorem}\label{additive_subset_product_of_scalar_multiples}
    Let $R$ be a commutative ring and $A$ be a subring of $R$.
    Let $a,b\in\carrier[R]$.
    Then $\additiveSubsetProduct{R}{
            \pointwiseProduct{\MultiplicativeMonoid{R}}{\{a\}}{A}}{
            \pointwiseProduct{\MultiplicativeMonoid{R}}{\{b\}}{A}}
        =\pointwiseProduct{\MultiplicativeMonoid{R}}{
            \{\ringMul[R](a,b)\}}{A}$.
\end{theorem}
\begin{proof}
    Let $M=\MultiplicativeMonoid{R}$.
    Let $X=\pointwiseProduct{M}{\{a\}}{A}$.
    Let $Y=\pointwiseProduct{M}{\{b\}}{A}$.
    Let $Z=\pointwiseProduct{M}{\{\ringMul[R](a,b)\}}{A}$.
    $R$ is a ring by \obv{commutative_ring_is_ring}.
    $M$ is a monoid by \obv{multiplicative_monoid_is_exact_monoid}.
    $\carrier[M]=\carrier[R]$ and $\mul[M]=\ringMul[R]$ by
        \auto{multiplicative_monoid_fields}.
    $A\subseteq\carrier[R]$ by \obv{subring_subseteq_carrier}.
    Hence $A\subseteq\carrier[M]$.
    $\{a\},\{b\}\subseteq\carrier[M]$ by
        \auto{subseteq,singleton_iff}.
    $\subgroup{A}{\AdditiveGroup{R}}$ by
        \obv{subring_additive_subgroup}.
    $\ringMul[R](a,b)\in\carrier[R]$ by \obv{ring_mul_type}.

    Show that $\pointwiseProduct{M}{X}{Y}=Z$.
    \begin{subproof}
        Show that every element of $\pointwiseProduct{M}{X}{Y}$ is an
            element of $Z$.
        \begin{subproof}
            Fix $z\in\pointwiseProduct{M}{X}{Y}$.
            Take $x,y$ such that $x\in X$ and $y\in Y$ and
                $z=\ringMul[R](x,y)$ by
                \auto{pointwise_product_iff}.
            Take $c\in A$ such that $x=\ringMul[R](a,c)$ by
                \auto{pointwise_product_iff,singleton_iff}.
            Take $d\in A$ such that $y=\ringMul[R](b,d)$ by
                \auto{pointwise_product_iff,singleton_iff}.
            $c,d\in\carrier[R]$ by \obv{subseteq}.
            $\ringMul[R](b,d),\ringMul[R](c,d),
                \ringMul[R](c,b),\ringMul[R](b,c)\in\carrier[R]$ by
                \obv{ring_mul_type}.
            $\ringMul[R](c,d)\in A$ by \obv{subring_mul}.
            $z=\ringMul[R](
                \ringMul[R](a,c),\ringMul[R](b,d))$ by \auto.
            $\ringMul[R](c,\ringMul[R](b,d))
                =\ringMul[R](\ringMul[R](c,b),d)$ by
                \obv{ring_mul_associative}.
            $\ringMul[R](c,b)=\ringMul[R](b,c)$ by
                \obv{commutative_ring_mul_commutative}.
            $\ringMul[R](\ringMul[R](b,c),d)
                =\ringMul[R](b,\ringMul[R](c,d))$ by
                \obv{ring_mul_associative}.
            Hence $\ringMul[R](c,\ringMul[R](b,d))
                =\ringMul[R](b,\ringMul[R](c,d))$ by \auto.
            $\ringMul[R](\ringMul[R](a,c),\ringMul[R](b,d))
                =\ringMul[R](a,\ringMul[R](c,\ringMul[R](b,d)))$ by
                \obv{ring_mul_associative}.
            Hence $z=\ringMul[R](a,
                \ringMul[R](c,\ringMul[R](b,d)))$ by \auto.
            Hence $z=\ringMul[R](a,
                \ringMul[R](b,\ringMul[R](c,d)))$ by \auto.
            $z=\ringMul[R](\ringMul[R](a,b),\ringMul[R](c,d))$ by
                \obv{ring_mul_associative}.
            Follows by
                \auto{pointwise_product_iff,singleton_iff}.
        \end{subproof}
        Show that every element of $Z$ is an element of
            $\pointwiseProduct{M}{X}{Y}$.
        \begin{subproof}
            Fix $z\in Z$.
            Take $c\in A$ such that
                $z=\ringMul[R](\ringMul[R](a,b),c)$ by
                \auto{pointwise_product_iff,singleton_iff}.
            $c\in\carrier[R]$ by \obv{subseteq}.
            $\ringOne[R]\in A$ by \obv{subring_one}.
            $\ringMul[R](a,\ringOne[R])\in X$ and
                $\ringMul[R](b,c)\in Y$ by
                \auto{pointwise_product_iff,singleton_iff}.
            $z=\ringMul[R](
                \ringMul[R](a,\ringOne[R]),\ringMul[R](b,c))$ by
                \auto{ring_mul_associative,ring_mul_one_right}.
            Follows by \auto{pointwise_product_iff}.
        \end{subproof}
        Follows by set extensionality.
    \end{subproof}
    $\subgroup{
        \pointwiseProduct{\MultiplicativeMonoid{R}}{
            \{\ringMul[R](a,b)\}}{A}}{\AdditiveGroup{R}}$ by
        \obv{left_scalar_multiple_is_additive_subgroup}.
    Hence $\subgroup{Z}{\AdditiveGroup{R}}$ by \auto.
    $\AdditiveGroup{R}$ is a group by \obv{additive_group_is_group}.
    $Z\subseteq\carrier[\AdditiveGroup{R}]$ by
        \obv{subgroup_subseteq}.
    $\generatedSubgroup{Z}{\AdditiveGroup{R}}=Z$ by
        \obv{generated_subgroup_of_subgroup}.
    Follows by \auto{additive_subset_product}.
\end{proof}

\begin{proposition}\label{additive_subset_product_singleton_scalar_multiple}
    Let $R$ be a ring and $A$ be a subring of $R$.
    Let $a,b\in\carrier[R]$.
    Then $\additiveSubsetProduct{R}{\{a\}}{
            \pointwiseProduct{\MultiplicativeMonoid{R}}{\{b\}}{A}}
        =\pointwiseProduct{\MultiplicativeMonoid{R}}{
            \{\ringMul[R](a,b)\}}{A}$.
\end{proposition}
\begin{proof}
    Let $M=\MultiplicativeMonoid{R}$.
    Let $H=\pointwiseProduct{M}{\{b\}}{A}$.
    $M$ is a monoid by \obv{multiplicative_monoid_is_exact_monoid}.
    Hence $M$ is a semigroup by \obv{monoid_implies_semigroup}.
    $\carrier[M]=\carrier[R]$ and $\mul[M]=\ringMul[R]$ by
        \auto{multiplicative_monoid_fields}.
    $A\subseteq\carrier[R]$ by \obv{subring_subseteq_carrier}.
    Hence $A\subseteq\carrier[M]$.
    $\{a\},\{b\}\subseteq\carrier[M]$ by
        \auto{subseteq,singleton_iff}.
    $\subgroup{A}{\AdditiveGroup{R}}$ by
        \obv{subring_additive_subgroup}.
    $\subgroup{
        \pointwiseProduct{\MultiplicativeMonoid{R}}{\{b\}}{A}}{
        \AdditiveGroup{R}}$ by
        \obv{left_scalar_multiple_is_additive_subgroup}.
    Hence $\subgroup{H}{\AdditiveGroup{R}}$ by \auto.
    $\additiveSubsetProduct{R}{\{a\}}{H}
        =\pointwiseProduct{\MultiplicativeMonoid{R}}{\{a\}}{H}$ by
        \obv{additive_subset_product_singleton_subgroup}.
    Hence $\additiveSubsetProduct{R}{\{a\}}{H}
        =\pointwiseProduct{M}{\{a\}}{H}$ by \auto.
    $\pointwiseProduct{M}{
            \pointwiseProduct{M}{\{a\}}{\{b\}}}{A}
        =\pointwiseProduct{M}{\{a\}}{
            \pointwiseProduct{M}{\{b\}}{A}}$ by
        \obv{pointwise_product_associative}.
    Hence $\pointwiseProduct{M}{\{a\}}{H}
        =\pointwiseProduct{M}{
            \pointwiseProduct{M}{\{a\}}{\{b\}}}{A}$ by
        \auto.
    $\pointwiseProduct{M}{\{a\}}{\{b\}}
        =\{\ringMul[R](a,b)\}$ by
        \auto{pointwise_product_singletons}.
    Hence $\additiveSubsetProduct{R}{\{a\}}{H}
        =\pointwiseProduct{M}{\{\ringMul[R](a,b)\}}{A}$ by
        \auto.
    Follows by \auto.
\end{proof}

\begin{proposition}\label{scalar_square_subring_subset_scalar_multiple}
    Let $R$ be a ring and $A$ be a subring of $R$.
    Let $a\in A$.
    Then $\pointwiseProduct{\MultiplicativeMonoid{R}}{
            \{\ringMul[R](a,a)\}}{A}
        \subseteq\pointwiseProduct{\MultiplicativeMonoid{R}}{\{a\}}{A}$.
\end{proposition}
\begin{proof}
    Let $M=\MultiplicativeMonoid{R}$.
    $\carrier[M]=\carrier[R]$ and $\mul[M]=\ringMul[R]$ by
        \auto{multiplicative_monoid_fields}.
    $A\subseteq\carrier[R]$ by \obv{subring_subseteq_carrier}.
    Hence $a\in\carrier[R]$.
    Show that every element of
        $\pointwiseProduct{M}{\{\ringMul[R](a,a)\}}{A}$ is an element of
        $\pointwiseProduct{M}{\{a\}}{A}$.
    \begin{subproof}
        Fix $z\in\pointwiseProduct{M}{\{\ringMul[R](a,a)\}}{A}$.
        Take $c\in A$ such that
            $z=\ringMul[R](\ringMul[R](a,a),c)$ by
            \auto{pointwise_product_iff,singleton_iff}.
        $c\in\carrier[R]$ by \obv{subseteq}.
        $\ringMul[R](a,c)\in A$ by \obv{subring_mul}.
        $\ringMul[R](\ringMul[R](a,a),c)
            =\ringMul[R](a,\ringMul[R](a,c))$ by
            \obv{ring_mul_associative}.
        Hence $z=\ringMul[R](a,\ringMul[R](a,c))$ by \auto.
        Follows by \auto{pointwise_product_iff,singleton_iff}.
    \end{subproof}
    Hence $\pointwiseProduct{M}{\{\ringMul[R](a,a)\}}{A}
        \subseteq\pointwiseProduct{M}{\{a\}}{A}$ by \obv{subseteq}.
    Follows by \auto.
\end{proof}

\section{The additive subset semigroup}

\begin{definition}\label{additive_subset_semigroup}
    $\AdditiveSubsetSemigroup{R}=\structure[magma]{
        \carrier=\pow{\carrier[R]},
        \mul=\Apply{\binaryGraphOnOp}{
            \pow{\carrier[R]},\pow{\carrier[R]},
            \fun{A,B:\Set}{\additiveSubsetProduct{R}{A}{B}}}
    }$.
\end{definition}

\begin{proposition}\label{additive_subset_semigroup_carrier}
    $\carrier[\AdditiveSubsetSemigroup{R}]=\pow{\carrier[R]}$.
\end{proposition}
\begin{proof}
    Follows by
        \auto{additive_subset_semigroup,magmaaggregateprojection_carrier}.
\end{proof}

\begin{proposition}\label{additive_subset_semigroup_multiplication}
    $\mul[\AdditiveSubsetSemigroup{R}]
        =\Apply{\binaryGraphOnOp}{
            \pow{\carrier[R]},\pow{\carrier[R]},
            \fun{A,B:\Set}{\additiveSubsetProduct{R}{A}{B}}}$.
\end{proposition}
\begin{proof}
    Follows by
        \auto{additive_subset_semigroup,magmaaggregateprojection_mul}.
\end{proof}

\begin{proposition}\label{additive_subset_semigroup_mul_apply}
    Suppose $A,B\subseteq\carrier[R]$.
    Then $\mul[\AdditiveSubsetSemigroup{R}](A,B)
        =\additiveSubsetProduct{R}{A}{B}$.
\end{proposition}
\begin{proof}
    $A,B\in\pow{\carrier[R]}$ by \obv{powerset_intro}.
    Follows by
        \auto{additive_subset_semigroup_multiplication,
            logical_binary_graph_on_apply}.
\end{proof}

\begin{theorem}\label{additive_subset_semigroup_is_exact}
    Let $R$ be a ring.
    Then $\AdditiveSubsetSemigroup{R}$ is an exact semigroup.
\end{theorem}
\begin{proof}
    Let $S=\AdditiveSubsetSemigroup{R}$.
    $\carrier[S]=\pow{\carrier[R]}$ by
        \auto{additive_subset_semigroup_carrier}.
    $\mul[S]=\Apply{\binaryGraphOnOp}{
        \pow{\carrier[R]},\pow{\carrier[R]},
        \fun{A,B:\Set}{\additiveSubsetProduct{R}{A}{B}}}$ by
        \auto{additive_subset_semigroup_multiplication}.
    Show that for all $A,B\in\pow{\carrier[R]}$ we have
        $\additiveSubsetProduct{R}{A}{B}\in\pow{\carrier[R]}$.
    \begin{subproof}
        Fix $A,B\in\pow{\carrier[R]}$.
        Hence $A,B\subseteq\carrier[R]$ by \obv{powerset_elim}.
        Hence $\additiveSubsetProduct{R}{A}{B}\subseteq\carrier[R]$ by
            \obv{additive_subset_product_type}.
        Follows by \obv{powerset_intro}.
    \end{subproof}
    Hence $\mul[S]\in
        \funs{\carrier[S]\times\carrier[S]}{\carrier[S]}$ by
        \auto{logical_binary_graph_on_function}.
    Show that for all $A,B,C\in\carrier[S]$ we have
        $\mul[S](A,\mul[S](B,C))
            =\mul[S](\mul[S](A,B),C)$.
    \begin{subproof}
        Fix $A,B,C\in\carrier[S]$.
        Hence $A,B,C\in\pow{\carrier[R]}$.
        Hence $A,B,C\subseteq\carrier[R]$ by \obv{powerset_elim}.
        $\additiveSubsetProduct{R}{A}{B}\subseteq\carrier[R]$ and
            $\additiveSubsetProduct{R}{B}{C}\subseteq\carrier[R]$ by
            \obv{additive_subset_product_type}.
        $\mul[S](A,B)=\additiveSubsetProduct{R}{A}{B}$ and
            $\mul[S](B,C)=\additiveSubsetProduct{R}{B}{C}$ by
            \auto{additive_subset_semigroup_mul_apply}.
        $\mul[S](A,\additiveSubsetProduct{R}{B}{C})
            =\additiveSubsetProduct{R}{A}{
                \additiveSubsetProduct{R}{B}{C}}$ and
        $\mul[S](\additiveSubsetProduct{R}{A}{B},C)
            =\additiveSubsetProduct{R}{
                \additiveSubsetProduct{R}{A}{B}}{C}$ by
            \auto{additive_subset_semigroup_mul_apply}.
        Follows by \auto{additive_subset_product_associative}.
    \end{subproof}
    Hence $S$ is a semigroup by
        \auto{additive_subset_semigroup,magmaaggregateexact,
            magma,semigroup}.
    Follows by
        \auto{additive_subset_semigroup,magmaaggregateexact}.
\end{proof}

\begin{classification}\label{additive_subset_semigroup_is_semigroup}
    Let $R$ be a ring.
    Then $\AdditiveSubsetSemigroup{R}$ is a semigroup.
\end{classification}
\begin{proof}
    $\AdditiveSubsetSemigroup{R}$ is an exact semigroup by
        \obv{additive_subset_semigroup_is_exact}.
\end{proof}

\section{Positive additively generated powers}

\begin{definition}\label{additive_subset_power}
    $\additiveSubsetPower{R}{U}{n}
        =\semigroupPositivePower{\AdditiveSubsetSemigroup{R}}{U}{n}$.
\end{definition}

\begin{proposition}\label{additive_subset_power_type}
    Let $R$ be a ring.
    Suppose $U\subseteq\carrier[R]$.
    Let $n$ be a natural number.
    Suppose $n\neq\emptyset$.
    Then $\additiveSubsetPower{R}{U}{n}\subseteq\carrier[R]$.
\end{proposition}
\begin{proof}
    Let $S=\AdditiveSubsetSemigroup{R}$.
    $S$ is a semigroup by \obv{additive_subset_semigroup_is_semigroup}.
    $\carrier[S]=\pow{\carrier[R]}$ by
        \auto{additive_subset_semigroup_carrier}.
    Hence $U\in\carrier[S]$ by \obv{powerset_intro}.
    $\semigroupPositivePower{S}{U}{n}\in\carrier[S]$ by
        \obv{semigroup_positive_power_type}.
    Follows by
        \auto{additive_subset_power,powerset_elim}.
\end{proof}

\begin{proposition}\label{additive_subset_power_one}
    Let $R$ be a ring.
    Suppose $U\subseteq\carrier[R]$.
    Then $\additiveSubsetPower{R}{U}{\suc{\emptyset}}=U$.
\end{proposition}
\begin{proof}
    Let $S=\AdditiveSubsetSemigroup{R}$.
    $S$ is a semigroup by \obv{additive_subset_semigroup_is_semigroup}.
    $\carrier[S]=\pow{\carrier[R]}$ by
        \auto{additive_subset_semigroup_carrier}.
    Hence $U\in\carrier[S]$ by \obv{powerset_intro}.
    Follows by
        \auto{additive_subset_power,semigroup_positive_power_one}.
\end{proof}

\begin{theorem}\label{pointwise_subset_power_subseteq_additive_subset_power}
    Let $R$ be a ring. Suppose $U\subseteq\carrier[R]$.
    Let $n$ be a natural number such that $n\neq\emptyset$.
    Then
        $\pointwiseSubsetPower{\MultiplicativeMonoid{R}}{U}{n}
            \subseteq\additiveSubsetPower{R}{U}{n}$.
\end{theorem}
\begin{proof}
    Let $M=\MultiplicativeMonoid{R}$.
    Let $S=\AdditiveSubsetSemigroup{R}$.
    $M$ is a monoid by \obv{multiplicative_monoid_is_exact_monoid}.
    Hence $M$ is a unital magma by \obv{monoid_is_unital_magma}.
    $S$ is a semigroup by \obv{additive_subset_semigroup_is_semigroup}.
    $\carrier[S]=\pow{\carrier[R]}$ by
        \auto{additive_subset_semigroup_carrier}.
    Hence $U\in\carrier[S]$ by \obv{powerset_intro}.
    $\carrier[M]=\carrier[R]$ and $\neutral[M]=\ringOne[R]$ by
        \auto{multiplicative_monoid_fields}.
    Hence $U\subseteq\carrier[M]$.
    Let $C=\{k\in\naturals\mid
        \pointwiseSubsetPower{M}{U}{\suc{k}}
            \subseteq\additiveSubsetPower{R}{U}{\suc{k}}\}$.

    $\emptyset$ is a natural number by \obv{omega_empty}.
    $\pointwiseSubsetPower{M}{U}{\emptyset}=\{\neutral[M]\}$ by
        \obv{pointwise_subset_power_zero}.
    $\pointwiseSubsetPower{M}{U}{\suc{\emptyset}}
        =\pointwiseProduct{M}{
            \pointwiseSubsetPower{M}{U}{\emptyset}}{U}$ by
        \obv{pointwise_subset_power_successor}.
    Hence $\pointwiseSubsetPower{M}{U}{\suc{\emptyset}}
        =\pointwiseProduct{M}{\{\neutral[M]\}}{U}$ by \auto.
    $\pointwiseProduct{M}{\{\neutral[M]\}}{U}=U$ by
        \obv{pointwise_product_neutral_singleton_left}.
    Hence $\pointwiseSubsetPower{M}{U}{\suc{\emptyset}}=U$ by
        \auto.
    $\additiveSubsetPower{R}{U}{\suc{\emptyset}}=U$ by
        \obv{additive_subset_power_one}.
    Hence $\emptyset\in C$ by \auto{omega_empty,subseteq_refl}.

    Show that for every $k\in C$ we have $\suc{k}\in C$.
    \begin{subproof}
        Fix $k\in C$.
        Hence $k\in\naturals$ and
            $\pointwiseSubsetPower{M}{U}{\suc{k}}
                \subseteq\additiveSubsetPower{R}{U}{\suc{k}}$.
        $\suc{k}$ is a natural number and $\suc{k}\neq\emptyset$ by
            \obv{omega_successor,suc_neq_emptyset}.
        $\pointwiseSubsetPower{M}{U}{\suc{k}}
            \subseteq\carrier[M]$ by
            \obv{pointwise_subset_power_type}.
        $\additiveSubsetPower{R}{U}{\suc{k}}
            \subseteq\carrier[R]$ by
            \obv{additive_subset_power_type}.
        Hence $\additiveSubsetPower{R}{U}{\suc{k}}
            \subseteq\carrier[M]$ by \auto.
        $\pointwiseSubsetPower{M}{U}{\suc{\suc{k}}}
            =\pointwiseProduct{M}{
                \pointwiseSubsetPower{M}{U}{\suc{k}}}{U}$ by
            \obv{pointwise_subset_power_successor}.
        $\semigroupPositivePower{S}{U}{\suc{\suc{k}}}
            =\mul[S](\semigroupPositivePower{S}{U}{\suc{k}},U)$ by
            \obv{semigroup_positive_power_successor}.
        $\semigroupPositivePower{S}{U}{\suc{k}}
            =\additiveSubsetPower{R}{U}{\suc{k}}$ by
            \auto{additive_subset_power}.
        $\mul[S](\additiveSubsetPower{R}{U}{\suc{k}},U)
            =\additiveSubsetProduct{R}{
                \additiveSubsetPower{R}{U}{\suc{k}}}{U}$ by
            \auto{additive_subset_semigroup_mul_apply}.
        Hence $\additiveSubsetPower{R}{U}{\suc{\suc{k}}}
            =\additiveSubsetProduct{R}{
                \additiveSubsetPower{R}{U}{\suc{k}}}{U}$ by
            \auto{additive_subset_power}.
        $\pointwiseProduct{M}{
                \pointwiseSubsetPower{M}{U}{\suc{k}}}{U}
            \subseteq\pointwiseProduct{M}{
                \additiveSubsetPower{R}{U}{\suc{k}}}{U}$ by
            \auto{pointwise_product_monotone,subseteq_refl}.
        $\pointwiseProduct{\MultiplicativeMonoid{R}}{
                \additiveSubsetPower{R}{U}{\suc{k}}}{U}
            \subseteq\additiveSubsetProduct{R}{
                \additiveSubsetPower{R}{U}{\suc{k}}}{U}$ by
            \obv{pointwise_product_subseteq_additive_subset_product}.
        Hence $\pointwiseProduct{M}{
                \additiveSubsetPower{R}{U}{\suc{k}}}{U}
            \subseteq\additiveSubsetProduct{R}{
                \additiveSubsetPower{R}{U}{\suc{k}}}{U}$ by
            \auto.
        Hence $\pointwiseSubsetPower{M}{U}{\suc{\suc{k}}}
            \subseteq\additiveSubsetPower{R}{U}{\suc{\suc{k}}}$ by
            \auto{subseteq_transitive}.
        $\suc{k}$ is a natural number by \obv{omega_successor}.
        Follows by \auto.
    \end{subproof}
    Hence $C$ is an inductive set by \obv{inductive_set}.
    Hence every natural number is an element of $C$ by
        \obv{omega_minimal}.
    $\unions{n}$ is a natural number and $n=\suc{\unions{n}}$ by
        \obv{natural_nonzero_successor_of_union}.
    Hence $\unions{n}\in C$.
    Follows by \auto.
\end{proof}

\begin{proposition}\label{additive_subset_power_successor}
    Let $R$ be a ring.
    Suppose $U\subseteq\carrier[R]$.
    Let $n$ be a natural number.
    Suppose $n\neq\emptyset$.
    Then $\additiveSubsetPower{R}{U}{\suc{n}}
        =\additiveSubsetProduct{R}{\additiveSubsetPower{R}{U}{n}}{U}$.
\end{proposition}
\begin{proof}
    Let $S=\AdditiveSubsetSemigroup{R}$.
    $S$ is a semigroup by \obv{additive_subset_semigroup_is_semigroup}.
    $\carrier[S]=\pow{\carrier[R]}$ by
        \auto{additive_subset_semigroup_carrier}.
    Hence $U\in\carrier[S]$ by \obv{powerset_intro}.
    $\semigroupPositivePower{S}{U}{n}\in\carrier[S]$ by
        \obv{semigroup_positive_power_type}.
    Hence $\semigroupPositivePower{S}{U}{n}\subseteq\carrier[R]$ by
        \auto{powerset_elim}.
    $\semigroupPositivePower{S}{U}{\suc{n}}
        =\mul[S](\semigroupPositivePower{S}{U}{n},U)$ by
        \obv{semigroup_positive_power_successor}.
    Follows by
        \auto{additive_subset_power,additive_subset_semigroup_mul_apply}.
\end{proof}

\begin{theorem}\label{additive_subset_power_of_scalar_multiple}
    Let $R$ be a commutative ring and $A$ be a subring of $R$.
    Let $a\in\carrier[R]$.
    Let $n$ be a natural number.
    Suppose $n\neq\emptyset$.
    Then $\additiveSubsetPower{R}{
            \pointwiseProduct{\MultiplicativeMonoid{R}}{\{a\}}{A}}{n}
        =\pointwiseProduct{\MultiplicativeMonoid{R}}{
            \{\ringPower{R}{a}{n}\}}{A}$.
\end{theorem}
\begin{proof}
    Let $M=\MultiplicativeMonoid{R}$.
    Let $U=\pointwiseProduct{M}{\{a\}}{A}$.
    $R$ is a ring by \obv{commutative_ring_is_ring}.
    $A\subseteq\carrier[R]$ by \obv{subring_subseteq_carrier}.
    $\subgroup{A}{\AdditiveGroup{R}}$ by
        \obv{subring_additive_subgroup}.
    $U\subseteq\carrier[R]$ by
        \auto{subgroup_subseteq,additive_group_fields,
            left_scalar_multiple_is_additive_subgroup}.
    Let $C=\{k\in\naturals\mid
        \additiveSubsetPower{R}{U}{\suc{k}}
            =\pointwiseProduct{M}{
                \{\ringPower{R}{a}{\suc{k}}\}}{A}\}$.
    $\additiveSubsetPower{R}{U}{\suc{\emptyset}}=U$ by
        \obv{additive_subset_power_one}.
    $\ringPower{R}{a}{\suc{\emptyset}}=a$ by
        \obv{ring_power_one}.
    Hence $\emptyset\in C$ by \auto{omega_empty}.

    Show that for every $k\in C$ we have $\suc{k}\in C$.
    \begin{subproof}
        Fix $k\in C$.
        Hence $k$ is a natural number and
            $\additiveSubsetPower{R}{U}{\suc{k}}
                =\pointwiseProduct{M}{
                    \{\ringPower{R}{a}{\suc{k}}\}}{A}$.
        $\suc{k}\in\naturals$ and $\suc{k}\neq\emptyset$ by
            \obv{omega_successor,suc_neq_emptyset}.
        $\ringPower{R}{a}{\suc{k}}\in\carrier[R]$ by
            \obv{ring_power_type}.
        $\additiveSubsetPower{R}{U}{\suc{\suc{k}}}
            =\additiveSubsetProduct{R}{
                \additiveSubsetPower{R}{U}{\suc{k}}}{U}$ by
            \obv{additive_subset_power_successor}.
        Hence $\additiveSubsetPower{R}{U}{\suc{\suc{k}}}
            =\pointwiseProduct{M}{
                \{\ringMul[R](\ringPower{R}{a}{\suc{k}},a)\}}{A}$ by
            \auto{additive_subset_product_of_scalar_multiples}.
        $\ringPower{R}{a}{\suc{\suc{k}}}
            =\ringMul[R](\ringPower{R}{a}{\suc{k}},a)$ by
            \obv{ring_power_successor}.
        Hence $\suc{k}\in C$ by \auto{omega_successor}.
    \end{subproof}
    Hence $C$ is an inductive set by \obv{inductive_set}.
    Hence every natural number is an element of $C$ by
        \obv{omega_minimal}.
    $\unions{n}$ is a natural number and $n=\suc{\unions{n}}$ by
        \obv{natural_nonzero_successor_of_union}.
    Hence $\unions{n}\in C$.
    Follows by \auto.
\end{proof}

\begin{theorem}\label{additive_subset_power_addition}
    Let $R$ be a ring.
    Suppose $U\subseteq\carrier[R]$.
    Let $m,n$ be natural numbers.
    Suppose $m\neq\emptyset$ and $n\neq\emptyset$.
    Then $\additiveSubsetPower{R}{U}{m+n}
        =\additiveSubsetProduct{R}{\additiveSubsetPower{R}{U}{m}}
            {\additiveSubsetPower{R}{U}{n}}$.
\end{theorem}
\begin{proof}
    Let $S=\AdditiveSubsetSemigroup{R}$.
    $S$ is a semigroup by \obv{additive_subset_semigroup_is_semigroup}.
    $\carrier[S]=\pow{\carrier[R]}$ by
        \auto{additive_subset_semigroup_carrier}.
    Hence $U\in\carrier[S]$ by \obv{powerset_intro}.
    $\semigroupPositivePower{S}{U}{m},
        \semigroupPositivePower{S}{U}{n}\in\carrier[S]$ by
        \obv{semigroup_positive_power_type}.
    Hence $\semigroupPositivePower{S}{U}{m},
        \semigroupPositivePower{S}{U}{n}\subseteq\carrier[R]$ by
        \auto{powerset_elim}.
    $\semigroupPositivePower{S}{U}{m+n}
        =\mul[S](\semigroupPositivePower{S}{U}{m},
            \semigroupPositivePower{S}{U}{n})$ by
        \obv{semigroup_positive_power_addition}.
    Follows by
        \auto{additive_subset_power,additive_subset_semigroup_mul_apply}.
\end{proof}

\subsection{Additive closure and contraction}

\begin{theorem}\label{additive_subset_power_is_subgroup}
    Let $R$ be a ring.
    Suppose $U\subseteq\carrier[R]$.
    Let $n$ be a natural number.
    Suppose $\suc{\suc{\emptyset}}\subseteq n$.
    Then
        $\subgroup{\additiveSubsetPower{R}{U}{n}}
            {\AdditiveGroup{R}}$.
\end{theorem}
\begin{proof}
    $\emptyset,\suc{\emptyset},\suc{\suc{\emptyset}}\in\naturals$ by
        \obv{omega_empty,omega_successor}.
    Take a natural number $k$ such that
        $n=\suc{\suc{\emptyset}}+k$ by
        \obv{natural_additive_difference}.
    $\emptyset+k=k$ by \obv{num_addition_zero_left}.
    $\suc{\emptyset}+k=\suc{k}$ and
        $\suc{\suc{\emptyset}}+k=\suc{\suc{k}}$ by
        \auto{num_addition_successor_left}.
    Hence $n=\suc{\suc{k}}$.
    $\suc{k}$ is a natural number and $\suc{k}\neq\emptyset$ by
        \obv{omega_successor,suc_neq_emptyset}.
    Hence $\additiveSubsetPower{R}{U}{\suc{k}}\subseteq\carrier[R]$ by
        \obv{additive_subset_power_type}.
    $\additiveSubsetPower{R}{U}{\suc{\suc{k}}}
        =\additiveSubsetProduct{R}{
            \additiveSubsetPower{R}{U}{\suc{k}}}{U}$ by
        \obv{additive_subset_power_successor}.
    $\subgroup{
        \additiveSubsetProduct{R}{
            \additiveSubsetPower{R}{U}{\suc{k}}}{U}}
        {\AdditiveGroup{R}}$ by
        \obv{additive_subset_product_is_subgroup}.
    Follows by \auto.
\end{proof}

\begin{theorem}\label{additive_subset_power_subset_of_square_subset}
    Let $R$ be a ring.
    Suppose $U\subseteq\carrier[R]$.
    Suppose $\additiveSubsetProduct{R}{U}{U}\subseteq U$.
    Let $n$ be a natural number.
    Suppose $n\neq\emptyset$.
    Then $\additiveSubsetPower{R}{U}{n}\subseteq U$.
\end{theorem}
\begin{proof}
    Let $C=\{k\in\naturals\mid
        \additiveSubsetPower{R}{U}{\suc{k}}\subseteq U\}$.
    $\additiveSubsetPower{R}{U}{\suc{\emptyset}}=U$ by
        \obv{additive_subset_power_one}.
    Hence $\emptyset\in C$ by
        \auto{omega_empty,subseteq_refl}.

    Show that for every $k\in C$ we have $\suc{k}\in C$.
    \begin{subproof}
        Fix $k\in C$.
        Hence $k$ is a natural number and
            $\additiveSubsetPower{R}{U}{\suc{k}}\subseteq U$.
        $\suc{k}$ is a natural number and $\suc{k}\neq\emptyset$ by
            \obv{omega_successor,suc_neq_emptyset}.
        Hence $\additiveSubsetPower{R}{U}{\suc{k}}
            \subseteq\carrier[R]$ by
            \obv{additive_subset_power_type}.
        $\additiveSubsetPower{R}{U}{\suc{\suc{k}}}
            =\additiveSubsetProduct{R}{
                \additiveSubsetPower{R}{U}{\suc{k}}}{U}$ by
            \obv{additive_subset_power_successor}.
        $\additiveSubsetProduct{R}{
            \additiveSubsetPower{R}{U}{\suc{k}}}{U}
            \subseteq\additiveSubsetProduct{R}{U}{U}$ by
            \auto{additive_subset_product_monotone,subseteq_refl}.
        Hence $\additiveSubsetPower{R}{U}{\suc{\suc{k}}}
            \subseteq\additiveSubsetProduct{R}{U}{U}$ by
            \auto.
        Hence $\additiveSubsetPower{R}{U}{\suc{\suc{k}}}
            \subseteq U$ by \obv{subseteq_transitive}.
        $\suc{k}$ is a natural number by \obv{omega_successor}.
        Follows by \auto.
    \end{subproof}
    Hence $C$ is an inductive set by \obv{inductive_set}.
    Hence every natural number is an element of $C$ by
        \obv{omega_minimal}.
    $\unions{n}$ is a natural number and $n=\suc{\unions{n}}$ by
        \obv{natural_nonzero_successor_of_union}.
    Hence $\unions{n}\in C$.
    Follows by \auto.
\end{proof}

\begin{theorem}\label{additive_subset_power_successor_subset}
    Let $R$ be a ring.
    Suppose $U\subseteq\carrier[R]$.
    Suppose $\additiveSubsetProduct{R}{U}{U}\subseteq U$.
    Let $n$ be a natural number.
    Suppose $n\neq\emptyset$.
    Then $\additiveSubsetPower{R}{U}{\suc{n}}
        \subseteq\additiveSubsetPower{R}{U}{n}$.
\end{theorem}
\begin{proof}
    Let $D=\{k\in\naturals\mid
        \additiveSubsetPower{R}{U}{\suc{\suc{k}}}
            \subseteq\additiveSubsetPower{R}{U}{\suc{k}}\}$.
    $\additiveSubsetPower{R}{U}{\suc{\emptyset}}=U$ by
        \obv{additive_subset_power_one}.
    $\suc{\emptyset}$ is a natural number and
        $\suc{\emptyset}\neq\emptyset$ by
        \obv{omega_empty,omega_successor,suc_neq_emptyset}.
    $\additiveSubsetPower{R}{U}{\suc{\suc{\emptyset}}}
        =\additiveSubsetProduct{R}{
            \additiveSubsetPower{R}{U}{\suc{\emptyset}}}{U}$ by
        \obv{additive_subset_power_successor}.
    Hence $\additiveSubsetPower{R}{U}{\suc{\suc{\emptyset}}}
        =\additiveSubsetProduct{R}{U}{U}$ by \auto.
    Hence $\emptyset\in D$ by \auto{omega_empty}.

    Show that for every $k\in D$ we have $\suc{k}\in D$.
    \begin{subproof}
        Fix $k\in D$.
        Hence $k$ is a natural number and
            $\additiveSubsetPower{R}{U}{\suc{\suc{k}}}
                \subseteq\additiveSubsetPower{R}{U}{\suc{k}}$.
        $\suc{k},\suc{\suc{k}}\in\naturals$ and
            $\suc{k}\neq\emptyset$ and
            $\suc{\suc{k}}\neq\emptyset$ by
            \obv{omega_successor,suc_neq_emptyset}.
        Hence $\additiveSubsetPower{R}{U}{\suc{k}},
            \additiveSubsetPower{R}{U}{\suc{\suc{k}}}
            \subseteq\carrier[R]$ by
            \obv{additive_subset_power_type}.
        $\additiveSubsetPower{R}{U}{\suc{\suc{k}}}
            =\additiveSubsetProduct{R}{
                \additiveSubsetPower{R}{U}{\suc{k}}}{U}$ and
        $\additiveSubsetPower{R}{U}{\suc{\suc{\suc{k}}}}
            =\additiveSubsetProduct{R}{
                \additiveSubsetPower{R}{U}{\suc{\suc{k}}}}{U}$ by
            \obv{additive_subset_power_successor}.
        $\additiveSubsetProduct{R}{
            \additiveSubsetPower{R}{U}{\suc{\suc{k}}}}{U}
            \subseteq\additiveSubsetProduct{R}{
                \additiveSubsetPower{R}{U}{\suc{k}}}{U}$ by
            \auto{additive_subset_product_monotone,subseteq_refl}.
        Hence $\additiveSubsetPower{R}{U}{\suc{\suc{\suc{k}}}}
            \subseteq\additiveSubsetPower{R}{U}{\suc{\suc{k}}}$ by
            \auto.
        $\suc{k}$ is a natural number by \obv{omega_successor}.
        Follows by \auto.
    \end{subproof}
    Hence $D$ is an inductive set by \obv{inductive_set}.
    Hence every natural number is an element of $D$ by
        \obv{omega_minimal}.
    $\unions{n}$ is a natural number and $n=\suc{\unions{n}}$ by
        \obv{natural_nonzero_successor_of_union}.
    Hence $\unions{n}\in D$.
    Follows by \auto.
\end{proof}

\begin{theorem}\label{additive_subset_power_antitone}
    Let $R$ be a ring.
    Suppose $U\subseteq\carrier[R]$.
    Suppose $\additiveSubsetProduct{R}{U}{U}\subseteq U$.
    Let $m,n$ be natural numbers.
    Suppose $m\neq\emptyset$ and $n\neq\emptyset$ and $m\subseteq n$.
    Then $\additiveSubsetPower{R}{U}{n}
        \subseteq\additiveSubsetPower{R}{U}{m}$.
\end{theorem}
\begin{proof}
    Take a natural number $d$ such that $n=m+d$ by
        \obv{natural_additive_difference}.
    Let $C=\{k\in\naturals\mid
        \additiveSubsetPower{R}{U}{m+k}
            \subseteq\additiveSubsetPower{R}{U}{m}\}$.
    $m+\emptyset=m$ by \obv{num_addition_zero_right}.
    Hence $\emptyset\in C$ by
        \auto{omega_empty,subseteq_refl}.

    Show that for every $k\in C$ we have $\suc{k}\in C$.
    \begin{subproof}
        Fix $k\in C$.
        Hence $k$ is a natural number and
            $\additiveSubsetPower{R}{U}{m+k}
                \subseteq\additiveSubsetPower{R}{U}{m}$.
        $m+k$ is a natural number by
            \obv{num_addition_on_naturals}.
        $\emptyset\in m$ by
            \obv{natural_nonzero_contains_zero}.
        $m\subseteq m+k$ by
            \obv{natural_addition_left_subseteq_sum}.
        Hence $\emptyset\in m+k$ by \obv{subseteq}.
        Hence $m+k\neq\emptyset$ by \auto.
        $\additiveSubsetPower{R}{U}{\suc{m+k}}
            \subseteq\additiveSubsetPower{R}{U}{m+k}$ by
            \obv{additive_subset_power_successor_subset}.
        $m+\suc{k}=\suc{m+k}$ by
            \obv{num_addition_successor_right}.
        Hence $\additiveSubsetPower{R}{U}{m+\suc{k}}
            \subseteq\additiveSubsetPower{R}{U}{m+k}$ by
            \auto.
        Hence $\additiveSubsetPower{R}{U}{m+\suc{k}}
            \subseteq\additiveSubsetPower{R}{U}{m}$ by
            \auto{subseteq_transitive}.
        $\suc{k}$ is a natural number by \obv{omega_successor}.
        Follows by \auto.
    \end{subproof}
    Hence $C$ is an inductive set by \obv{inductive_set}.
    Hence every natural number is an element of $C$ by
        \obv{omega_minimal}.
    Hence $d\in C$.
    Follows by \auto.
\end{proof}

\begin{corollary}\label{pointwise_square_contraction_implies_additive_square_contraction}
    Let $R$ be a ring.
    Suppose $U\subseteq\carrier[R]$.
    Suppose $\subgroup{U}{\AdditiveGroup{R}}$.
    Suppose
        $\pointwiseProduct{\MultiplicativeMonoid{R}}{U}{U}\subseteq U$.
    Then $\additiveSubsetProduct{R}{U}{U}\subseteq U$.
\end{corollary}
\begin{proof}
    Let $G=\AdditiveGroup{R}$ and $M=\MultiplicativeMonoid{R}$.
    $\carrier[G]=\carrier[R]$ and $\carrier[M]=\carrier[R]$ by
        \auto{additive_group_fields,multiplicative_monoid_fields}.
    Hence $U\subseteq\carrier[M]$.
    Hence $\pointwiseProduct{M}{U}{U}\subseteq\carrier[M]$ by
        \obv{pointwise_product_type}.
    Hence $\pointwiseProduct{M}{U}{U}\subseteq\carrier[G]$.
    $\subgroup{U}{G}$ and
        $\pointwiseProduct{M}{U}{U}\subseteq U$ by \auto.
    $\generatedSubgroup{\pointwiseProduct{M}{U}{U}}{G}\subseteq U$ by
        \obv{generated_subgroup_least}.
    Follows by \auto{additive_subset_product}.
\end{proof}
